\chapter{Week 8 - Fourier Series}

\section{Introduction}

\subsection{Motivation}

\begin{itemize}
    \item \textbf{Taylor series:} Provides a local representation of a function $f: \R \to \R$ that is sufficiently regular ($C^\infty$) as an infinite sum of monomials:
    \[
    f(x) = \sum_{n=0}^\infty a_n x^n, \quad \text{where } a_n = \frac{f^{(n)}(0)}{n!} \quad \text{around } x = 0.
    \]
    For instance, $f(x) = \mathrm{e}^x = \sum_{n=0}^\infty \frac{x^n}{n!}$ for all $x \in \R$.
    
    \item \textbf{Fourier series:} Approximates and represents a periodic function by an infinite superposition of basic harmonic oscillations (sines and cosines). Given a $T$-periodic function $f: \R \to \R$, we seek an expansion of the form:
    \[
    f(x) \sim \frac{a_0}{2} + \sum_{n=1}^\infty \left[ a_n \cos\!\left(\frac{2\pi n}{T} x\right) + b_n \sin\!\left(\frac{2\pi n}{T} x\right) \right], \quad x \in \R.
    \]
\end{itemize}

\subsection{Preliminary Definitions and Results}

\begin{definition}{Periodicity}{periodicity}
Let $T > 0$. A function $f: \R \to \R$ is said to be \textbf{$T$-periodic} if:
\[
f(x + T) = f(x), \quad \forall x \in \R.
\]
\end{definition}

\begin{exemple}[Fundamental Harmonics]
For any $n \in \N^* = \{1, 2, 3, \dots\}$, the trigonometric functions:
\[
u_n(x) = \cos\!\left(\frac{2\pi n}{T} x\right) \quad \text{and} \quad v_n(x) = \sin\!\left(\frac{2\pi n}{T} x\right)
\]
are periodic of period $\frac{T}{n}$, and therefore are also $T$-periodic. Indeed:
\[
u_n\!\left(x + \frac{T}{n}\right) = \cos\!\left(\frac{2\pi n}{T}\left(x + \frac{T}{n}\right)\right) = \cos\!\left(\frac{2\pi n}{T} x + 2\pi\right) = \cos\!\left(\frac{2\pi n}{T} x\right) = u_n(x).
\]
\end{exemple}

\begin{definition}{Piecewise Continuous Function}{piecewise_continuous}
A function $f: \R \to \R$ is \textbf{piecewise continuous} if:
\begin{enumerate}
    \item It has at most a finite number of discontinuities on every bounded interval $[a, b] \subset \R$.
    \item At every point of discontinuity $x \in \R$, both the right-hand and left-hand limits exist and are finite:
    \[
    f(x+0) := \lim_{\substack{t \to x \\ t > x}} f(t) \quad \text{and} \quad f(x-0) := \lim_{\substack{t \to x \\ t < x}} f(t).
    \]
\end{enumerate}
\end{definition}

\begin{definition}{Piecewise $C^1$ Function}{piecewise_c1}
A function $f: \R \to \R$ is said to be \textbf{piecewise $C^1$} if:
\begin{enumerate}
    \item $f$ is piecewise continuous on $\R$.
    \item $f$ is continuously differentiable on each open subinterval between consecutive discontinuity points.
    \item At every discontinuity point $x \in \R$, the one-sided derivative limits exist and are finite:
    \[
    f'(x+0) := \lim_{\substack{t \to x \\ t > x}} f'(t) \quad \text{and} \quad f'(x-0) := \lim_{\substack{t \to x \\ t < x}} f'(t).
    \]
\end{enumerate}
\end{definition}

\begin{proposition}{Orthogonality Relations of Trigonometric Functions}{trig_orthogonality}
Let $m, n \in \N^*$ and $T > 0$. The trigonometric system satisfies the following orthogonality identities:
\begin{align}
\frac{2}{T} \int_0^T \cos\!\left(\frac{2\pi n}{T} x\right) \cos\!\left(\frac{2\pi m}{T} x\right) \dx &= \begin{cases} 0 & \text{if } m \neq n, \\ 1 & \text{if } m = n, \end{cases} \label{eq:ortho_cos} \\
\frac{2}{T} \int_0^T \sin\!\left(\frac{2\pi n}{T} x\right) \sin\!\left(\frac{2\pi m}{T} x\right) \dx &= \begin{cases} 0 & \text{if } m \neq n, \\ 1 & \text{if } m = n, \end{cases} \label{eq:ortho_sin} \\
\frac{2}{T} \int_0^T \sin\!\left(\frac{2\pi n}{T} x\right) \cos\!\left(\frac{2\pi m}{T} x\right) \dx &= 0, \quad \forall m, n \in \N^*. \label{eq:ortho_sincos}
\end{align}
\end{proposition}

\begin{demonstration}
Without loss of generality, assume $T = 2\pi$ via the change of variable $\theta = \frac{2\pi}{T} x$.
\begin{itemize}
    \item \textbf{Identity \eqref{eq:ortho_cos}:} Using the product-to-sum formula $\cos(a)\cos(b) = \frac{1}{2}[\cos(a-b) + \cos(a+b)]$:
    \[
    I = \frac{1}{\pi} \int_0^{2\pi} \cos(nx)\cos(mx) \dx = \frac{1}{2\pi} \int_0^{2\pi} \cos((n-m)x) \dx + \frac{1}{2\pi} \int_0^{2\pi} \cos((n+m)x) \dx.
    \]
    If $m \neq n$, since $n-m \neq 0$ and $n+m \neq 0$:
    \[
    I = \frac{1}{2\pi} \left[ \frac{\sin((n-m)x)}{n-m} \right]_0^{2\pi} + \frac{1}{2\pi} \left[ \frac{\sin((n+m)x)}{n+m} \right]_0^{2\pi} = 0 + 0 = 0.
    \]
    If $m = n$:
    \[
    I = \frac{1}{2\pi} \int_0^{2\pi} 1 \dx + \frac{1}{2\pi} \left[ \frac{\sin(2nx)}{2n} \right]_0^{2\pi} = \frac{2\pi}{2\pi} + 0 = 1.
    \]
    \item \textbf{Identity \eqref{eq:ortho_sincos}:} Using $\sin(a)\cos(b) = \frac{1}{2}[\sin(a-b) + \sin(a+b)]$:
    \[
    I = \frac{1}{\pi} \int_0^{2\pi} \sin(nx)\cos(mx) \dx = \frac{1}{2\pi} \int_0^{2\pi} \sin((n-m)x) \dx + \frac{1}{2\pi} \int_0^{2\pi} \sin((n+m)x) \dx.
    \]
    If $m \neq n$:
    \[
    I = \frac{1}{2\pi} \left[ -\frac{\cos((n-m)x)}{n-m} \right]_0^{2\pi} + \frac{1}{2\pi} \left[ -\frac{\cos((n+m)x)}{n+m} \right]_0^{2\pi} = 0 + 0 = 0.
    \]
    If $m = n$, $\sin((n-m)x) = \sin(0) = 0$, giving again $I = 0$.
\end{itemize}
Identity \eqref{eq:ortho_sin} is established identically using $\sin(a)\sin(b) = \frac{1}{2}[\cos(a-b) - \cos(a+b)]$.
\end{demonstration}

\begin{lemme}{Integration of Periodic Functions}{periodic_integral}
Let $f: \R \to \R$ be a $T$-periodic, piecewise continuous function. Then the integral of $f$ over any interval of length $T$ is invariant:
\[
\int_a^{a+T} f(x) \dx = \int_0^T f(x) \dx, \quad \forall a \in \R.
\]
\end{lemme}

\begin{demonstration}
Split the integral at the boundary points $0$ and $T$:
\[
\int_a^{a+T} f(x) \dx = \int_a^0 f(x) \dx + \int_0^T f(x) \dx + \int_T^{a+T} f(x) \dx.
\]
In the third integral, substitute $x = y + T$ so that $\dd x = \dd y$. Since $f$ is $T$-periodic, $f(y+T) = f(y)$:
\[
\int_T^{a+T} f(x) \dx = \int_0^a f(y+T) \dd y = \int_0^a f(y) \dd y.
\]
Consequently:
\[
\int_a^{a+T} f(x) \dx = -\int_0^a f(x) \dx + \int_0^T f(x) \dx + \int_0^a f(y) \dd y = \int_0^T f(x) \dx.
\]
\end{demonstration}

\section{Fourier Series of a Periodic Function}

\subsection{Definition and Fourier Coefficients}

\begin{definition}{Fourier Series and Partial Sums}{fourier_series}
Let $f: \R \to \R$ be a $T$-periodic, piecewise continuous function.
\begin{enumerate}
    \item For any integer $N \ge 1$, the \textbf{Fourier partial sum of order $N$} of $f$ is defined by:
    \[
    \mathcal{F}_N f(x) = \frac{a_0}{2} + \sum_{n=1}^N \left[ a_n \cos\!\left(\frac{2\pi n}{T} x\right) + b_n \sin\!\left(\frac{2\pi n}{T} x\right) \right],
    \]
    where the real scalars $a_n$ ($n \ge 0$) and $b_n$ ($n \ge 1$) are the \textbf{Fourier coefficients} of $f$, defined by:
    \begin{align*}
    a_n &= \frac{2}{T} \int_0^T f(x) \cos\!\left(\frac{2\pi n}{T} x\right) \dx \quad (n \ge 0), \\
    b_n &= \frac{2}{T} \int_0^T f(x) \sin\!\left(\frac{2\pi n}{T} x\right) \dx \quad (n \ge 1).
    \end{align*}
    \item The \textbf{Fourier series} of $f$ is the formal pointwise limit as $N \to \infty$:
    \[
    \mathcal{F}f(x) := \lim_{N \to \infty} \mathcal{F}_N f(x) = \frac{a_0}{2} + \sum_{n=1}^\infty \left[ a_n \cos\!\left(\frac{2\pi n}{T} x\right) + b_n \sin\!\left(\frac{2\pi n}{T} x\right) \right].
    \]
\end{enumerate}
\end{definition}

\begin{remarque}[Choice of the Constant Term $\frac{a_0}{2}$]
Setting $n = 0$ in the formula for $a_n$ gives $a_0 = \frac{2}{T} \int_0^T f(x) \dx$. The constant term in the expansion is therefore:
\[
\frac{a_0}{2} = \frac{1}{T} \int_0^T f(x) \dx,
\]
which represents exactly the mean value of $f$ over one period.
\end{remarque}

\subsection{Heuristic Justification: Hilbert Space Projection}

To understand the algebraic structure underlying Fourier coefficients, set $T = 2\pi$ for simplicity. Let $V$ denote the real vector space of continuous $2\pi$-periodic functions $g: \R \to \R$.

Equip $V$ with the $L^2$-inner product:
\[
\inner{f}{g} := \frac{1}{\pi} \int_0^{2\pi} f(x) g(x) \dx.
\]
By the orthogonality relations (Proposition \ref{prop:trig_orthogonality}):
\begin{align*}
\inner{\cos(nx)}{\cos(mx)} &= \delta_{nm}, \\
\inner{\sin(nx)}{\sin(mx)} &= \delta_{nm}, \\
\inner{\cos(nx)}{\sin(mx)} &= 0.
\end{align*}
For the constant basis function $u_0(x) \equiv 1$:
\[
\inner{1}{1} = \frac{1}{\pi} \int_0^{2\pi} 1 \dx = 2, \quad \inner{\cos(nx)}{1} = 0, \quad \inner{\sin(nx)}{1} = 0.
\]
Assuming the trigonometric family $\{1, \cos(nx), \sin(nx)\}_{n \ge 1}$ forms an orthogonal basis of $V$, any function $f \in V$ decomposes into its orthogonal projection components:
\[
f(x) = c \cdot 1 + \sum_{n=1}^\infty \left[ a_n \cos(nx) + b_n \sin(nx) \right],
\]
where the coefficients are given by:
\begin{align*}
c &= \frac{\inner{f}{1}}{\inner{1}{1}} = \frac{\frac{1}{\pi} \int_0^{2\pi} f(x) \dx}{2} = \frac{a_0}{2}, \\
a_n &= \frac{\inner{f}{\cos(nx)}}{\inner{\cos(nx)}{\cos(nx)}} = \frac{1}{\pi} \int_0^{2\pi} f(x) \cos(nx) \dx, \\
b_n &= \frac{\inner{f}{\sin(nx)}}{\inner{\sin(nx)}{\sin(nx)}} = \frac{1}{\pi} \int_0^{2\pi} f(x) \sin(nx) \dx.
\end{align*}

\subsection{Convergence of Fourier Series}

\begin{theoreme}{Dirichlet's Theorem}{dirichlet}
Let $f: \R \to \R$ be a $T$-periodic, \textbf{piecewise $C^1$} function. Then for every $x \in \R$, the Fourier series converges pointwisely:
\[
\mathcal{F}f(x) = \lim_{N \to \infty} \mathcal{F}_N f(x) = \frac{f(x+0) + f(x-0)}{2}.
\]
In particular, at every point $x \in \R$ where $f$ is continuous:
\[
\mathcal{F}f(x) = f(x).
\]
\end{theoreme}

\begin{remarque}[Physical Interpretation]
Dirichlet's theorem confirms that any piecewise $C^1$ signal can be synthesized exactly as a countable superposition of pure harmonics, with the Fourier series converging to the midpoint of the jump discontinuity at every point of discontinuity.
\end{remarque}

\subsection{Worked Examples}

\begin{exemple}[Square Wave]
Consider the $2\pi$-periodic function defined on $[0, 2\pi)$ by:
\[
f(x) = \begin{cases} 1 & \text{if } x \in [0, \pi), \\ 0 & \text{if } x \in [\pi, 2\pi). \end{cases}
\]

\textbf{1. Computation of the Fourier coefficients:}
\begin{align*}
a_0 &= \frac{1}{\pi} \int_0^{2\pi} f(x) \dx = \frac{1}{\pi} \int_0^\pi 1 \dx + \frac{1}{\pi} \int_\pi^{2\pi} 0 \dx = 1.
\end{align*}
For $n \ge 1$:
\begin{align*}
a_n &= \frac{1}{\pi} \int_0^{2\pi} f(x) \cos(nx) \dx = \frac{1}{\pi} \int_0^\pi \cos(nx) \dx = \frac{1}{\pi} \left[ \frac{\sin(nx)}{n} \right]_0^\pi = 0. \\
b_n &= \frac{1}{\pi} \int_0^{2\pi} f(x) \sin(nx) \dx = \frac{1}{\pi} \int_0^\pi \sin(nx) \dx = \frac{1}{\pi} \left[ -\frac{\cos(nx)}{n} \right]_0^\pi \\
    &= -\frac{1}{\pi n} \left( \cos(n\pi) - 1 \right) = \frac{1 - (-1)^n}{\pi n} = \begin{cases} 0 & \text{if } n \text{ is even}, \\ \dfrac{2}{\pi n} & \text{if } n \text{ is odd}. \end{cases}
\end{align*}

\textbf{2. Fourier series expansion:}
Writing $n = 2k + 1$ with $k \ge 0$:
\[
\mathcal{F}f(x) = \frac{1}{2} + \frac{2}{\pi} \sum_{k=0}^\infty \frac{\sin((2k+1)x)}{2k+1}.
\]

\textbf{3. Dirichlet's theorem evaluation:}
\begin{itemize}
    \item On the open intervals $x \in (0, \pi)$ and $x \in (\pi, 2\pi)$, $f$ is continuous, hence $\mathcal{F}f(x) = f(x)$.
    \item At the discontinuity points $x = 0$ and $x = \pi$:
    \[
    \mathcal{F}f(0) = \frac{f(0+0) + f(0-0)}{2} = \frac{1 + 0}{2} = \frac{1}{2},
    \]
    \[
    \mathcal{F}f(\pi) = \frac{f(\pi+0) + f(\pi-0)}{2} = \frac{0 + 1}{2} = \frac{1}{2}.
    \]
\end{itemize}
\end{exemple}

\begin{exemple}[Sawtooth Wave]
Consider the $2\pi$-periodic function defined on $[0, 2\pi)$ by:
\[
f(x) = x, \quad x \in [0, 2\pi).
\]

\textbf{1. Computation of the Fourier coefficients:}
\[
a_0 = \frac{1}{\pi} \int_0^{2\pi} x \dx = \frac{1}{\pi} \left[ \frac{x^2}{2} \right]_0^{2\pi} = \frac{4\pi^2}{2\pi} = 2\pi.
\]
For $n \ge 1$, we integrate by parts using $u = x \implies \dd u = \dd x$, and $\dd v = \cos(nx)\dx \implies v = \frac{\sin(nx)}{n}$:
\begin{align*}
a_n &= \frac{1}{\pi} \int_0^{2\pi} x \cos(nx) \dx = \frac{1}{\pi} \left( \left[ \frac{x \sin(nx)}{n} \right]_0^{2\pi} - \int_0^{2\pi} \frac{\sin(nx)}{n} \dx \right) \\
    &= \frac{1}{\pi} \left( 0 + \left[ \frac{\cos(nx)}{n^2} \right]_0^{2\pi} \right) = \frac{1}{\pi n^2} (1 - 1) = 0.
\end{align*}
For $b_n$, integrate by parts with $u = x \implies \dd u = \dd x$, and $\dd v = \sin(nx)\dx \implies v = -\frac{\cos(nx)}{n}$:
\begin{align*}
b_n &= \frac{1}{\pi} \int_0^{2\pi} x \sin(nx) \dx = \frac{1}{\pi} \left( \left[ -\frac{x \cos(nx)}{n} \right]_0^{2\pi} + \int_0^{2\pi} \frac{\cos(nx)}{n} \dx \right) \\
    &= \frac{1}{\pi} \left( -\frac{2\pi \cos(2\pi n)}{n} + 0 + \left[ \frac{\sin(nx)}{n^2} \right]_0^{2\pi} \right) = -\frac{2}{n}.
\end{align*}

\textbf{2. Fourier series expansion:}
\[
\mathcal{F}f(x) = \frac{a_0}{2} + \sum_{n=1}^\infty b_n \sin(nx) = \pi - 2 \sum_{n=1}^\infty \frac{\sin(nx)}{n}.
\]

\textbf{3. Dirichlet's theorem evaluation:}
\begin{itemize}
    \item For $x \in (0, 2\pi)$, $f$ is continuous: $\mathcal{F}f(x) = x$.
    \item At $x = 0$ (and $x = 2k\pi$):
    \[
    \mathcal{F}f(0) = \frac{f(0+0) + f(0-0)}{2} = \frac{0 + 2\pi}{2} = \pi.
    \]
\end{itemize}
\end{exemple}